% Generated draft: requires mathematical and source review.
\documentclass[11pt]{article}
\usepackage[margin=25mm]{geometry}
\usepackage{fontspec}
\setmainfont{Latin Modern Roman}
\usepackage{amsmath,amssymb,mathtools,unicode-math}
\setmathfont{Latin Modern Math}
\usepackage{enumitem,needspace,xurl,hyperref}
\setlength{\parindent}{0pt}
\setlength{\parskip}{5pt}
\setlength{\emergencystretch}{3em}
\newcommand{\N}{\mathbb{N}}
\newcommand{\Z}{\mathbb{Z}}
\newcommand{\Q}{\mathbb{Q}}
\newcommand{\R}{\mathbb{R}}
\newcommand{\C}{\mathbb{C}}
\newcommand{\F}{\mathbb{F}}
\newcommand{\A}{\mathbb{A}}
\newcommand{\PP}{\mathbb{P}}
\newcommand{\E}{\mathbb{E}}
\newcommand{\eps}{\varepsilon}
\newcommand{\dd}{\,\mathrm{d}}
\begin{document}
\section*{UnsolvedMath ID: 5500035}
\textbf{Source problem code:} catalogue entry
\subsection{Definitions and mathematical statement}
Let $p_0,\ldots,p_{n-1}$ be the initial robot locations in a metric space $(X,d)$, with only the robot at the designated source $p_0$ initially awake. Under the usual unit-speed convention, an awake robot can move between locations in time given by distance, and awakens a sleeping robot when they meet. The newly awakened robot may move and awaken others immediately. A schedule's makespan is the time when all robots are awake; write $t^*$ for the minimum achievable makespan.

The first question is whether computing an optimal schedule is NP-hard when the sites lie in $\R^2$ with Euclidean distance, or with $L_1$ distance $d(x,y)=|x_1-y_1|+|x_2-y_2|$. For a formal complexity claim, finite input encoding and the optimization or threshold-decision formulation must be fixed; NP-hardness does not assert membership in NP for exact Euclidean comparisons.

The second question is whether general finite metrics admit a polynomial-time approximation with a guarantee asymptotically better than $O(\log n)$: a schedule of length at most $\rho(n)t^*$ with, for example, $\rho(n)=o(\log n)$. The source does not state an exact improvement scale; a smaller constant in a logarithmic factor is a different target.

The source also describes an equivalent rooted spanning-tree problem, with maximum out-degree two and radius equal to the maximum sum of edge lengths along a root-to-vertex path. A scheduling convention is missing in that description: with only one initially awake robot, the source initially has only one active continuation, whereas awakening an ordinary vertex creates two. A usual awakening-tree representation therefore needs a root constraint or auxiliary-vertex convention, beyond an unrestricted binary branching root. The source's tree assertion is preserved with this qualification rather than used to redefine the physical scheduling problem.
\subsection{Short English statement}
Determine the planar hardness of optimally waking robots, and improve the general-metric logarithmic approximation ratio. The supplied report says planar hardness is settled while the approximation question remains separate.
\subsection{Sources}
\begin{itemize}
\item Local UnsolvedMath registry, permanent ID 5500035, exact source code catalogue entry, ``Freeze-Tag: Optimal Strategies for Awakening a Swarm of Robots''; no proposer or year supplied. Registry status solved, classification SOLVED-IN-LITERATURE.
\item \emph{The Open Problems Project}, Problem 35, \url{https://topp.openproblem.net/p35}; public source specification \url{https://github.com/edemaine/topp}. Supplied source-TeX extraction accessed July 29, 2026. The maintained entry is reported as conjecturing Euclidean/$L_1$ planar NP-hardness and reporting NP-completeness for star metrics.
\item The research summary says planar NP-hardness is established in Euclidean and $L_p$ planes, citing Yu et al. (2017) and a strong-NP-hardness result (2025), without precise publications or norm/parameter hypotheses. It explicitly leaves the general-metric improvement beyond the $O(\log n)$ approximation scale as a separate open question. These are supplied assessments, not independent literature verification.
\item Rights note: the source repository LICENSE expressly covers software; reuse permission for problem text is flagged NEEDS\_REVIEW. Public availability does not settle that scope.
\end{itemize}
\end{document}
